%! Diagonalizing a Matrix — Unit 6, Lesson 6.2 — Guided Notes (Answer Key)
\documentclass[10pt]{article}
\usepackage{linalg-article}
\usepackage{linalg-key}   % pulls in -boxes; do NOT also load -boxes
\newcommand{\TallMath}[1]{$\displaystyle #1\rule[-1.4em]{0pt}{3.2em}$}
\newcommand{\xx}{\mathbf{x}}
\newcommand{\zero}{\mathbf{0}}

\begin{document}

\pageheader{Unit 6, Lesson 6.2}{Guided Notes --- Answer Key}
\namedateperiod

\vspace{0.12in}

\begin{objectivebox}
\textbf{By the end of this lesson, I will be able to\dots}
\begin{itemize}\setlength{\itemsep}{2pt}
  \item line eigenvectors into the columns of $X$ and eigenvalues down the diagonal of $\Lambda$, and explain why $AX=X\Lambda$;
  \item \emph{diagonalize} a matrix as $A=X\Lambda X^{-1}$ when it has enough independent eigenvectors;
  \item use $A^k=X\Lambda^k X^{-1}$ to compute powers quickly, and say when a matrix \emph{cannot} be diagonalized.
\end{itemize}
\end{objectivebox}

\vspace{0.06in}

\begin{vocabbox}
\small
Fill in each term as we build it together.
\vspace{2pt}
\vocabans{Eigenvector matrix $X$ (eigenvectors in the columns)}{put each eigenvector of $A$ into a column of $X$ (order matches $\Lambda$)}
\vocabans{Eigenvalue matrix $\Lambda$ (eigenvalues down the diagonal)}{a diagonal matrix with the eigenvalues on the diagonal, in the same order as $X$'s columns}
\vocabans{Diagonalization $A=X\Lambda X^{-1}$}{rebuild $A$ from its eigenvectors and eigenvalues; possible when $X$ is invertible}
\vocabans{Power rule $A^k=X\Lambda^k X^{-1}$}{powers of $A$ reduce to powers of the eigenvalues (the inner $X^{-1}X$'s cancel)}
\end{vocabbox}

\begin{hookbox}
Yesterday you found that $A=\begin{bsmallmatrix} 2 & 1\\ 1 & 2 \end{bsmallmatrix}$ has eigenvalues
$\lambda=3,1$ with eigenvectors $\begin{bsmallmatrix}1\\1\end{bsmallmatrix}$ and $\begin{bsmallmatrix}1\\-1\end{bsmallmatrix}$.
\begin{itemize}\setlength{\itemsep}{1pt}
  \item Suppose you needed $A^{10}$. Would you really multiply $A$ by itself ten times? \ansline{No --- there is a shortcut: raise each eigenvalue to the 10th power ($3^{10}$ and $1^{10}$) and rebuild.}
  \item Along an eigenvector, $A$ acts like a single \emph{number} $\lambda$. Could we make the whole matrix act that simply? \ansline{Yes --- in eigenvector coordinates $A$ becomes the diagonal matrix $\Lambda$.}
\end{itemize}
\end{hookbox}

\begin{notesbox}{1. Line Up the Eigenvectors: $AX=X\Lambda$}
Put the eigenvectors of $A$ into the \textbf{columns} of a matrix $X$, and the eigenvalues down the
diagonal of $\Lambda$ (same order):
\[
  X=\begin{bmatrix} 1 & 1\\ 1 & -1 \end{bmatrix}\ \text{(eigenvectors)},\qquad
  \Lambda=\begin{bmatrix} 3 & 0\\ 0 & 1 \end{bmatrix}\ \text{(eigenvalues)}.
\]
Now multiply $A$ by $X$ \emph{one column at a time}. Each column is an eigenvector, so $A$ just scales it:
\[
  AX=\big[\,A\xx_1\ \ A\xx_2\,\big]=\big[\,3\xx_1\ \ 1\xx_2\,\big]
  =\begin{bmatrix} 1 & 1\\ 1 & -1 \end{bmatrix}\begin{bmatrix} 3 & 0\\ 0 & 1 \end{bmatrix}=X\Lambda.
\]
So lining the eigenvectors up gives the clean identity $AX={\color{keyred}\mathbf{X\Lambda}}$. (Scaling the columns of $X$
by the eigenvalues is exactly what multiplying by $\Lambda$ on the \ans{right} does.)
\end{notesbox}

\begin{notesbox}{2. Solve for $A$: the Diagonalization $A=X\Lambda X^{-1}$}
If the eigenvectors are \textbf{independent}, then $X$ is invertible. Multiply $AX=X\Lambda$ on the right
by $X^{-1}$:
\[
  AX X^{-1}=X\Lambda X^{-1}\quad\Longrightarrow\quad A={\color{keyred}\mathbf{X\Lambda X^{-1}}}.
\]
This is the \textbf{diagonalization} of $A$. (Multiplying instead by $X^{-1}$ on the \emph{left} gives
$\Lambda=X^{-1}AX$ --- the matrix is \emph{diagonal} in eigenvector coordinates.)

\vspace{3pt}
For our matrix, $X=\begin{bsmallmatrix}1&1\\1&-1\end{bsmallmatrix}$ has $\det X=-2$, so by the Unit~2
formula $X^{-1}=\frac{1}{-2}\begin{bsmallmatrix}-1&-1\\-1&1\end{bsmallmatrix}=\begin{bsmallmatrix}\frac12&\frac12\\[2pt]\frac12&-\frac12\end{bsmallmatrix}$. Then
\[
  X\Lambda X^{-1}
  =\begin{bmatrix} 3 & 1\\ 3 & -1 \end{bmatrix}\begin{bmatrix} \tfrac12 & \tfrac12\\[2pt] \tfrac12 & -\tfrac12 \end{bmatrix}
  =\begin{bmatrix}{\color{keyred}\mathbf{2}} & {\color{keyred}\mathbf{1}}\\[3pt]{\color{keyred}\mathbf{1}} & {\color{keyred}\mathbf{2}}\end{bmatrix}=A.\quad\checkmark
\]
\end{notesbox}

\begin{notesbox}{3. The Payoff: Powers Become Easy}
Watch what happens to $A^2$ when we write $A=X\Lambda X^{-1}$:
\[
  A^2=(X\Lambda X^{-1})(X\Lambda X^{-1})=X\Lambda \underbrace{(X^{-1}X)}_{I}\Lambda X^{-1}
  =X\Lambda^2 X^{-1}.
\]
The inside $X^{-1}X$ collapses to $I$. The same cancellation repeats, so for \emph{any} power
\[
  A^k=X\,\Lambda^k\,X^{-1},\qquad\text{and }\Lambda^k=\begin{bmatrix}\lambda_1^k & 0\\ 0 & \lambda_2^k\end{bmatrix}
  \text{ is trivial (just powers of the }\lambda\text{'s).}
\]
Try $A^2$ for our matrix: $\Lambda^2=\begin{bsmallmatrix}9&0\\0&1\end{bsmallmatrix}$, so
$A^2=X\Lambda^2 X^{-1}=\begin{bsmallmatrix}{\color{keyred}\mathbf{5}}&{\color{keyred}\mathbf{4}}\\[2pt]{\color{keyred}\mathbf{4}}&{\color{keyred}\mathbf{5}}\end{bsmallmatrix}$.
(Check by squaring $A$ directly!)

\begin{center}
\begin{tikzpicture}[scale=1, font=\scriptsize,
    box/.style={draw=charcoal, rounded corners=1pt, fill=blush, minimum height=8mm, inner sep=3pt, align=center},
    ar/.style={->, thick, burgundy}]
  \node[box] (v)  at (0,0)   {any vector\\ $\mathbf{v}$};
  \node[box] (c)  at (3.3,0) {eigenvector\\ coordinates};
  \node[box] (sc) at (7.1,0) {scaled\\ coordinates};
  \node[box] (av) at (10.4,0){result\\ $A\mathbf{v}$};
  \draw[ar] (v)  -- node[above]{$X^{-1}$} node[below]{\itshape rewrite} (c);
  \draw[ar] (c)  -- node[above]{$\Lambda$} node[below]{\itshape scale by $\lambda$'s} (sc);
  \draw[ar] (sc) -- node[above]{$X$} node[below]{\itshape rebuild} (av);
\end{tikzpicture}
\end{center}
$A=X\Lambda X^{-1}$ reads right to left: \textbf{rewrite} $\mathbf v$ in eigenvector coordinates, \textbf{scale}
each by its $\lambda$, then \textbf{rebuild}. All of $A$'s work is the easy middle step.
\end{notesbox}

\begin{notesbox}{4. When Can We Diagonalize? + The Road Ahead}
The one requirement: $A$ needs enough \textbf{independent eigenvectors} to fill the columns of an
invertible $X$ (for a $2\times2$, two independent eigenvectors).
\begin{itemize}\setlength{\itemsep}{2pt}
  \item \textbf{Different eigenvalues $\Rightarrow$ always works.} Distinct eigenvalues have independent eigenvectors.
  \item \textbf{A repeat can fail.} The shear $\begin{bsmallmatrix}2&1\\0&2\end{bsmallmatrix}$ has only $\lambda=2$
        and \emph{one} eigenvector direction $\begin{bsmallmatrix}1\\0\end{bsmallmatrix}$ --- not enough to build $X$, so it \ans{cannot} be diagonalized.
\end{itemize}

\vspace{3pt}
\textbf{The road ahead.} \textbf{6.3} shows \emph{symmetric} matrices always diagonalize --- and with
\emph{perpendicular} eigenvectors --- and \textbf{6.4} uses $A^k=X\Lambda^k X^{-1}$ to run a real system
forward in time.
\end{notesbox}

\begin{practicebox}
Show each step.
\begin{enumerate}[leftmargin=1.6em, itemsep=6pt]
  \item $A=\begin{bsmallmatrix} 3 & 1\\ 0 & 2 \end{bsmallmatrix}$ has $\lambda=3,2$ with eigenvectors $\begin{bsmallmatrix}1\\0\end{bsmallmatrix}$ and $\begin{bsmallmatrix}1\\-1\end{bsmallmatrix}$. Write $X$, $\Lambda$, and $X^{-1}$. \ansline{$X=\begin{bsmallmatrix}1&1\\0&-1\end{bsmallmatrix}$, $\Lambda=\begin{bsmallmatrix}3&0\\0&2\end{bsmallmatrix}$; $\det X=-1$, so $X^{-1}=\begin{bsmallmatrix}1&1\\0&-1\end{bsmallmatrix}$.}
  \item Using part 1, write $A^2$ as $X\Lambda^2 X^{-1}$ and multiply it out. \ansline{$\Lambda^2=\begin{bsmallmatrix}9&0\\0&4\end{bsmallmatrix}$, so $A^2=X\Lambda^2 X^{-1}=\begin{bsmallmatrix}9&4\\0&-4\end{bsmallmatrix}\begin{bsmallmatrix}1&1\\0&-1\end{bsmallmatrix}=\begin{bsmallmatrix}9&5\\0&4\end{bsmallmatrix}$. (Check: $A^2=\begin{bsmallmatrix}3&1\\0&2\end{bsmallmatrix}^2=\begin{bsmallmatrix}9&5\\0&4\end{bsmallmatrix}$.)}
  \item Why does the shear $\begin{bsmallmatrix}2&1\\0&2\end{bsmallmatrix}$ fail to diagonalize? \ansline{It has a repeated eigenvalue $\lambda=2$ with only one eigenvector direction $\begin{bsmallmatrix}1\\0\end{bsmallmatrix}$, so $X$ would have dependent columns and no inverse.}
\end{enumerate}
\end{practicebox}

\begin{teachernote}
The spine is a four-step build on \emph{yesterday's} matrix $\begin{bsmallmatrix}2&1\\1&2\end{bsmallmatrix}$
($\lambda=3,1$; eigenvectors $\begin{bsmallmatrix}1\\1\end{bsmallmatrix},\begin{bsmallmatrix}1\\-1\end{bsmallmatrix}$),
so nothing new has to be computed --- we \emph{reuse} their answer. §1: line the eigenvectors into $X$ and
watch $AX=X\Lambda$ column by column. §2: right-multiply by $X^{-1}$ to get $A=X\Lambda X^{-1}$ (the Unit~2
inverse formula gives the halves). §3: the payoff --- $X^{-1}X$ cancels, so $A^k=X\Lambda^k X^{-1}$; the
pipeline figure reads $A$ right-to-left as rewrite $\to$ scale $\to$ rebuild. §4: diagonalizable iff enough
independent eigenvectors; the shear is the standard counterexample. Common errors: putting eigenvalues in an
order that doesn't match the columns of $X$; forgetting $\Lambda$ multiplies \emph{columns} (on the right);
and fraction slips in $X^{-1}$. Emphasize that the practice matrix (triangular, $\det X=-1$) is the
\emph{fraction-free} version to build confidence.
\end{teachernote}

\end{document}
